\documentclass{article}
\usepackage[margin=1in]{geometry}
\usepackage{hyperref}
\usepackage{parskip}
\usepackage{lmodern}

\pagestyle{empty}

\begin{document}

\begin{flushleft}
    {\Large Alexandre Ait Ettajer} \\
    \href{mailto:alexandre.aitettajer@gmail.com}{alexandre.aitettajer@gmail.com} \quad +1(281) 435-9481 \quad Boston, MA \\
    GitHub: \href{https://github.com/aiteta}{Aiteta} \quad LinkedIn: \href{https://www.linkedin.com/in/alexandre-ait-ettajer-65740b170/}{alexandre-ait-ettajer} \\
\end{flushleft}

\vspace{1em}

Dear Open Source Server team,

I'm applying for the Software Engineer II role on the Open Source Server team. I'm a Go programmer with four years building production backend systems at MathWorks and an M.S. from UT Austin in distributed and parallel systems. I want to go deeper on distributed systems, and I'd rather do it on the core of a system people actually run than anywhere else.

The most relevant thing I've built is a fault-tolerant distributed file system in Go, using Paxos for replication across nodes. That forced me to work through the exact concepts your posting names: fault tolerance, consistency tradeoffs, replication, and what failure handling looks like in practice rather than on paper.

At MathWorks I rebuilt the client-server messaging behind MATLAB visualizations, steered a cross-team working group on distributed systems architecture, and did the operational work too: resolving mass build and test failures, running cross-cluster pipelines, shipping performance improvements. I also like ambiguous starting points. Prototyping MATLAB's first hierarchical visualization meant turning a vague idea into something shippable, which sounds like the day-to-day of your team.

I'm excited to learn a large, serious codebase and contribute to the reliability and operability of the server itself.

Best, \\
Alexandre Ait Ettajer

\end{document}
